\documentclass[11pt,letterpaper]{article}
\usepackage[T1]{fontenc}
\usepackage{lmodern,microtype,amsmath,amssymb,amsthm,booktabs,array,geometry,xcolor,hyperref,fancyhdr}
\geometry{left=1.0in,right=1.0in,top=.88in,bottom=.89in,headheight=14pt}
\definecolor{ink}{HTML}{19384A}\definecolor{blue}{HTML}{126F70}
\hypersetup{colorlinks=true,urlcolor=blue,linkcolor=blue,citecolor=blue,pdftitle={Sparse equations for singular moduli}}
\pagestyle{fancy}\fancyhf{}\fancyhead[L]{\small\sffamily\color{ink}Sparse equations for singular moduli}
\fancyhead[R]{\small\sffamily\color{ink}Research note}\fancyfoot[R]{\sffamily\small\thepage}
\setlength{\parindent}{0pt}\setlength{\parskip}{.52em}
\newtheorem{theorem}{Theorem}\newtheorem{lemma}[theorem]{Lemma}
\newtheorem{corollary}[theorem]{Corollary}
\newcommand{\Q}{\mathbb Q}\newcommand{\Z}{\mathbb Z}\newcommand{\F}{\mathbb F}
\newcommand{\supp}{\operatorname{supp}}\newcommand{\ord}{\operatorname{ord}}
\newcommand{\repo}[2]{\href{https://github.com/lsernau/grh-lean/blob/72fead102f663d5c7a407964a13dc0cf9d977207/grh/#1}{#2}}
\newcommand{\src}[1]{\textsf{\small #1}}\newcommand{\divider}{\par\noindent\textcolor{blue}{\rule{\linewidth}{.8pt}}\par}
\newcommand{\heading}[1]{\textbf{\color{ink}#1}}
\begin{document}
\thispagestyle{plain}
{\small\sffamily\color{blue} NUMBER THEORY \quad / \quad RESEARCH NOTE \hfill 25 NOVEMBER 2026}\par
\vspace{.35cm}
{\LARGE\bfseries\color{ink} Sparse equations for singular moduli}\par
{\large A local determinant argument after the Lean proof of Dirichlet GRH}
\divider
Let $j$ be the modular invariant and let $\alpha=j(\tau)$, where $\tau$ is imaginary quadratic. Such an $\alpha$ is a \emph{singular modulus}. Write $\Delta=-D<0$ for the discriminant of its endomorphism order, $H_\Delta$ for its monic class polynomial and $h=h(\Delta)=\deg H_\Delta$ for the class number. Define
\[
 s(\Delta)=\min\{\#\supp P : 0\ne P\in\Q[X],\ P(\alpha)=0\}.
\]
All conjugates have the same minimum. The degree of $P$ is unrestricted. The class polynomial gives $s(\Delta)\le h+1$.

\begin{theorem}[Four-term classification]\label{th:four}
A singular modulus is a root of a nonzero rational polynomial with at most four terms if and only if its degree over $\Q$ is at most three. In particular, every quartic singular modulus has minimum support exactly five.
\end{theorem}
Riffaut's conjecture concerns \emph{three} terms. Bilu, Luca and Pizarro-Madariaga proved it under Dirichlet GRH \cite{riffaut}. The upstream Lean theorem supplies that hypothesis. For four terms, an incoming ProveIt manuscript \cite{prior} proves an explicit cutoff and reduces the problem to 50 quartic and 29 quintic discriminants. We close all 79 using a $p$-adic argument uniform in the exponents. A stronger unconditional reduction leaves just five possible discriminants for a five-term equation in degree at least five: $-87,-131,-348,-432,-675$ (Theorem~\ref{th:five}). All five exclude equations of width at most $65{,}536$ (Corollary~\ref{cor:width}).

\begin{theorem}[Growth of the minimum support]\label{th:growth}
Uniformly over imaginary quadratic orders, the upstream Dirichlet theorem gives
\[
 s(-D)\gg \frac{h(-D)}{(\log(2D))^2},
 \qquad \lim_{D\to\infty}\frac{\log s(-D)}{\log D}=\tfrac12 .
\]
Independently of RH, $s(-D)\gg_\varepsilon D^{1/2-1/(4\sqrt e)-\varepsilon}$ for every $\varepsilon>0$. On a set of order discriminants of relative density $9/13$, the stronger unconditional estimate $s(-D)=h(-D)+1-O(\log(2h(-D)))$ holds.
\end{theorem}
For the family in Theorem~\ref{th:extreme}, the minimum equals $h+1$ exactly, including $1{,}048{,}577$ at $\Delta=-2^{44}$. The incoming lower bound $\gg D^{1/4}/\sqrt{\log D}$ \cite{prior} was improved in a parallel report to $\gg D^{1/4}/\sqrt{\log\log D}$ under GRH and $\gg_\varepsilon D^{1/4-\varepsilon}$ unconditionally \cite{parallel}. Theorem~\ref{th:growth} determines the exponent. The new results are mathematical proofs with exact arithmetic certificates. The upstream Lean verification covers the zero-location statements.
\clearpage
\section{Which idea transfers from the Lean development?}
The readable edition of \repo{src/Main.lean}{\texttt{Main.lean}} exports Mathlib's RH statement, its explicit zeta-zero formulation and three formulations for Dirichlet characters. In particular, \texttt{LFunction\_ne\_zero\_of\_half\_lt\_re} covers all complex Dirichlet characters at $\Re s>1/2$, with the principal character at $s=1$ excepted. \repo{src/Statements.lean}{The definitions of the statements} import Mathlib. Andrew Sutherland reports independent builds and kernel replays of both editions and agreement on the standard three axioms. We inspected source at commit \texttt{72fead10} and used the published verification record \cite{lean,audit}.

The comparison in \repo{src/New/Reduction/MainInequality.lean}{\texttt{New/Reduction/MainInequality}} concerns one corrected covolume. Short integral forms control its size and index from above. A spanning statement gives enough forms. A lower bound for determinants uses auxiliary twists chosen place by place and a global height estimate. The common main term cancels in \texttt{star\_inequality}. The more abstract \repo{src/New/Reduction/Lattices/TwistedLattice.lean}{\texttt{covolume\_add\_index\_nonneg}} still requires explicit operators, determinant-height control and vectors completing a basis. The resulting inequality bounds a weighted prime remainder in \repo{src/New/Reduction/OneSidedBound.lean}{\texttt{OneSidedBound}}. Its auxiliary modulus and quadratic twist are fixed before the scale tends to infinity. A Mellin transform records the difference between zero multiplicities for the original and twisted characters. A twist chosen to avoid cancellation and \repo{src/Lib/Analytic/Landau.lean}{Landau\textquotesingle s theorem} then give zero-freeness \cite{lean}.

Here the place-by-place determinant test has a different arithmetic realization. A $t$-term relation forces all $t$-row minors of $(\alpha_i^{n_j})$ to vanish, where $\alpha_i$ are conjugates and $n_j$ are the exponents. The incoming report used archimedean sizes to prove a minor nonzero when some conjugates have widely separated moduli. We select a finite place where many conjugates are \emph{close} after scaling. On a sufficiently small multiplicative disk, a sparse polynomial has few roots, uniformly in its exponents. This certifies full rank even for conjugates with almost equal or equal complex absolute values. The growth arguments select such places across an entire family of orders.

The local root-count argument and its CM applications are developed below. The Lean covolume theorem itself has hypotheses about unitary groups, theta lifts and auxiliary forms. Our application draws on the choice of a useful finite place and comparison of ranks, together with the exported Dirichlet theorem for the analytic estimates. It requires no zero theorem for general Hecke or Artin $L$-functions.

\heading{Relation to earlier work.} Other incoming packets already study class-group exponents, fields of coefficients, nonabelian prime-race criteria, and small splitting primes. The particular starting point here is the four-term question and the explicit residual list in \src{grh\_cm\_prime\_applications.zip}. Its analytic cutoff and discriminant sieves are credited below. The local fewnomial literature includes sharper results in various ranges, especially Lenstra and Avenda\~no--Krick \cite{few}. A parallel workspace manuscript \cite{parallel} also closes the 79 cases, with an independent finite sieve and a proof free of GRH. The large-discriminant part of Section~\ref{sec:five} extends its split/nonsplit-prime argument. Our stronger asymptotic estimates, the exact infinite family, the five-term reduction, and the constraints at $-131$ go beyond the results stated there. A literature search identified the published trinomial theorem; priority for the extensions is undetermined.
\clearpage
\section{A root count uniform in the exponents}\label{sec:local}
Normalize $v=v_p$ by $v(p)=1$ in finite extensions of $\Q_p$. A multiplicative disk of relative depth $r$ about $a\ne0$ is given by $v(x/a-1)\ge r$.
\begin{lemma}\label{lem:disk}
Let $P(X)=\sum_{j=1}^t c_jX^{n_j}$, with distinct nonnegative integer exponents and nonzero coefficients in a finite extension of $\Q_p$. Suppose $r$ is rational and $r>1/(p-1)$. Put $w_r(k)=rk-v_p(k!)$. In a disk of relative depth $r$, the number of distinct roots is at most
\[
 B_p(t,r)=\max\{k\ge0 : w_r(k)\le\max_{0\le i<t} w_r(i)\}.
\]
The bound counts algebraic roots in the disk.
\end{lemma}
\begin{proof}
Enlarge the finite extension to contain the roots, center and an element $b$ of valuation $r$ (rational depths suffice throughout). The substitutions $x=a\exp(by)$ and $y=b^{-1}\log(x/a)$ identify this disk with $v(y)\ge0$. In the resulting analytic function the coefficient of $y^k$ is $b^k M_k/k!$, where $M_k=\sum_j c_j a^{n_j}n_j^k$; the zeroth power is one. The first $t$ moments are not all zero, by the ordinary Vandermonde determinant in the distinct integers $n_j$. All later moments obey the monic recurrence with characteristic polynomial $\prod_j(T-n_j)\in\Z[T]$. Thus every $v(M_k)$ is at least $m=\min_{i<t}v(M_i)$. Some $i<t$ attains $m$. When $w_r(k)>\max_{i<t}w_r(i)$ the coefficient of index $k$ has strictly larger valuation than that of this $i$. Also $w_r(k)\to\infty$ by Legendre's factorial formula. Strassmann's theorem bounds the number of zeros by the last index of a coefficient of largest absolute value. This proves the claim.
\end{proof}
For $t=4,5$ it suffices that the relative depths of $t$ distinct roots are all strictly greater than
\[
 \kappa(p,t)=
 \begin{cases}2,&p=2,\\1,&p=3,\\1/(p-t+1),&p\ge5.\end{cases}
 \tag{1}
\]
Indeed, for $i<t$ and $n\ge t$ we need $(v_p(n!)-v_p(i!))/(n-i)\le\kappa$. If $p\ge t$, the numerator vanishes for $n<p$ and is at most $(n-1)/(p-1)$ otherwise. For $n\ge p$, dividing by $n-i$ gives at most $1/(p-i)$ when $i\ge1$; for $i=0$ the bound $1/(p-1)$ suffices. Each is at most $1/(p-t+1)$. For $p=2,3$ and $n\ge7$, use $(n-1)/(p-1)\le\kappa(n-4)$. Direct factorials for $4\le n\le6$ verify the remaining pairs $i<t$, $n\ge t$. For fewer terms the same bound follows from the monotonicity in $t$ of Lemma~\ref{lem:disk}. This argument allows rational depths in ramified fields and also disks centered at a nonunit after a common scaling.
\clearpage
\section{The earlier four-term reduction}\label{sec:prior}
For a primitive reduced quadratic form $[a,b,c]$ of discriminant $-D$, the corresponding conjugate has modulus within $2079$ of $\exp(\pi\sqrt D/a)$ \cite{riffaut}. This gives an archimedean complement to the disk lemma. If $t$ conjugates of increasing moduli $r_i$ have successive ratios at least 16, every generalized Vandermonde minor on them and $t$ distinct exponents $n_1<\cdots<n_t$ is nonzero. Scale column $i+1$ relative to column $i$ by $(r_ir_{i+1})^{-(n_{i+1}-n_i)/2}$. Moving away from the diagonal reduces entries by at least a factor four at every step. Each row's off-diagonal sum is at most $2\sum_{k\ge1}4^{-k}=2/3$ times its diagonal magnitude. Strict diagonal dominance proves the assertion.

For consecutive leading coefficients $a<b$ the integer inequalities
\[
       D\ge16a^2,\qquad 9D(b-a)^2\ge16a^2b^2                         \tag{2}
\]
suffice for a ratio greater than 16. Indeed the first exponential is at least $e^{4\pi}>34\cdot2079$ and the exponential ratio is at least $e^{4\pi/3}>32$, so the two cusp errors preserve the claimed separation.

\heading{What was inherited and verified.} The incoming CM report \cite{prior} applied the quadratic factorization $\zeta_K=\zeta L(\cdot,\chi_{d_K})$ and Greni\'e--Molteni's Corollary 1.2 \cite{gm} for $D\ge10^{11}$. It obtained an odd split prime coprime to the conductor in each of $(0,.046R]$, $(.048R,.189R]$, $(.230R,.500R]$ for $R=\sqrt D$. With the principal form these give four separated conjugates. We reproduced its exact endpoint checks, yielding margins $68,191,751$. Its finite sieves gave the following counts; we reran both.
\begin{center}\small
\begin{tabular}{@{}p{.26\linewidth}p{.65\linewidth}@{}}\toprule
Range or step & Outcome \\ \midrule
$10^7\le D<10^{11}$ &
Of $49{,}995{,}000{,}000$ order discriminants, $12{,}498{,}750{,}000$ pass at 2, $37{,}449{,}133{,}804$ via powers of an odd split prime and $47{,}116{,}196$ via three individually checked primes. \\[2pt]
$3\le D\le10^7$ &
Of $5{,}000{,}000$ discriminants, $67$ have $h<4$. Integer separation certifies $4{,}999{,}603$. Of the remaining $330$, interval separation removes $228$ and a determinant comparison removes $22$. \\[2pt]
Residual $80$ & $50$ have degree four, $29$ degree five, and $-4075$ has degree six. The earlier report covers $-4075$ by complex intervals. \\ \bottomrule
\end{tabular}
\end{center}
The disk proof gives more at $-4075$. At 7 all six conjugates are in one nonzero residue disk of depth at least one, by Section~\ref{sec:growth}. The proof of (1) for $p\ge t$ also works for $t=6$, giving threshold $1/2$ at 7. Hence $s(-4075)=7$. The following sections close the 79 earlier cases and provide a separate unconditional global proof through the stronger five-term reduction.

\clearpage
\section{Closing the four-term classification}\label{sec:cert}
The class polynomial $H(X)=\sum A_iX^i$ has distinct algebraic integer roots. For integers $v\ge0$ and $r$ prime to $p$, form
\[
 H(p^v(r+Z))=\sum_i B_i Z^i,\qquad
 B_i=\sum_{j\ge i} A_jp^{vj}\binom ji r^{j-i}.
\]
For any rational $u$, the \emph{least} minimizing index of $v_p(B_i)+iu$ equals the number of roots $z$ with $v_p(z)>u$, with multiplicity. To see this, choose $\pi$ with $v_p(\pi)=u$, scale $Z=\pi Y$, and normalize the polynomial by a coefficient of smallest valuation. Factoring into its linear factors in a splitting field and reducing shows that each root deeper than $u$ contributes a factor $Y$; each root at depth $u$ contributes $Y$ minus a nonzero residue, and each shallower root contributes a nonzero constant. This proves the assertion, including ties. For the original roots $\alpha$, $v_p(\alpha/(p^vr)-1)=v_p(z)$.

Set $u=\kappa(p,t)$. A least minimizing index at least $t$ excludes every relation of at most $t$ terms by Lemma~\ref{lem:disk}. Choose a rational depth strictly between $u$ and the shallowest of the selected depths. The eight triples $(p,v,r)$ used in all 79 certificates are
\[
 (2,3,5),\ (2,6,3),\ (2,6,5),\ (3,3,1),\ (3,3,2),\ (7,0,6),\ (13,0,5),\ (17,0,8).
\]

\heading{An example one can check by hand.} For $\Delta=-55$, the valuations of the five coefficients of $H(27(1+Z))$, constant term first, are $(20,18,17,14,12)$. With $p=3$, $t=4$ and $u=1$ the weights are $(20,19,19,17,16)$. The unique minimum is at index 4. All four roots are therefore in the required disk, so $s(-55)=5$. For $\Delta=-131$ the analogous valuations at $(17,0,8)$ are $(4,3,2,1,0,0)$. With $u=1/14$ the unique minimum is at index 4, enough for the four-term theorem.

\heading{Exact input and outcome.} The supplied generator enumerates primitive reduced forms and computes $j=1728E_4^3/(E_4^3-E_6^2)$ with rational interval endpoints. It encloses $\pi$ by Machin's formula, the exponential by Taylor remainders and repeated squaring, and $E_4,E_6$ by at least 160 terms with a tail bound from $|q|\le1/200$ and $\sigma_k(n)\le2n^k$ for $k=3,5$. Every coefficient interval of the product contains exactly one integer. The small checker then performs only integer and rational operations. All 50 quartics pass with $t=4$. Of the 29 quintics, 28 pass with $t=5$ and $-131$ passes with $t=4$. The other direction in Theorem~\ref{th:four} follows from $\#\supp H\le h+1$ for $h\le3$. Independently, the class-number formula and the fundamental-discriminant bounds of Watkins \cite{watkins} enumerate all 84 quartic and 29 quintic orders. The archive certifies each quartic directly. Indeed $h(d_K)\mid h$ by class group surjectivity, $\varphi(f)\le3h/h(d_K)$ and $\varphi(f)^2\ge f/2$ prime-powerwise, so the conductor enumeration is finite. The quintic list is also the complete class-number-five order list \cite{orders}. Thus all but possibly one quintic have $s=6$.

All these conclusions also hold over the CM field $K$. Ring class theory gives $[K(\alpha):K]=h=[\Q(\alpha):\Q]$. The powers of $\alpha$ retain their rational linear independence after extending scalars to $K$.
\clearpage
\section{Five terms leave at most five discriminants}\label{sec:five}
\begin{theorem}\label{th:five}
Every singular modulus of degree at most four satisfies an equation with at most five nonzero rational coefficients. Among higher degrees, such an equation can occur only at the order discriminants
\[
              -87,\quad -131,\quad -348,\quad -432,\quad -675.
\]
The degree is five at $-131$ and six at the other four. In particular, every singular modulus of degree at least seven requires at least six terms. These assertions use no hypothesis about zeros.
\end{theorem}
The low-degree assertion follows from $s\le h+1$. The certificates in Section~\ref{sec:cert} handle $h=5$. We now treat $h\ge6$. Write $\Delta=f^2d_K=-D$ and put $A=4\cdot2401^2=23{,}059{,}204$.

\heading{Nonsplit primes.} For $p\equiv1\pmod3$, $p\nmid f$ nonsplit in $K$, all conjugates reduce to nonzero supersingular residues. The ring class field is unramified over $K$ at $p$, so the absolute ramification index is at most two even when $p\mid d_K$. Conjugates with a common residue have relative depth at least $1/2$. If their number is at least five and $p\ge7$, (1) applies since $1/(p-4)<1/2$. The numbers of supersingular residues for $p=7,13,19,31$ are respectively $1,1,2,3$. In general for primes $p\equiv1\pmod3$, the number $S_p$ equals $(p-1)/12$ if $p\equiv1\pmod{12}$ and $(p+5)/12$ otherwise \cite{deuring}. Thus $h>4S_p$ suffices.

\heading{Split and conductor primes.} If $p\nmid f$ splits in $K$, every $a=p^j$ with $4a^2\le D$ occurs as a leading coefficient: lift a nonzero root of $b^2\equiv-D\pmod p$ to $p^j$, choose the parity of $b$, and take $|b|\le a$. Then $[a,b,(b^2+D)/(4a)]$ is integral, primitive and reduced, up to the boundary sign convention. For an odd prime $p\mid f$, the coefficient $a=p^2$ is available when $D\ge4p^4$. Take $b=pv$ of suitable parity and avoid the at most two residues for which $v^2+D/p^2\equiv0\pmod p$. Choose $|v|\le p$ by CRT. This makes $c=(v^2+D/p^2)/4$ coprime to $p$ and $c\ge p^2$.

For $D\ge A$, try $7,13,19,31$ in order until a prime does not divide $f$. If it is nonsplit, the preceding disk test applies. In the last two rows involving a prime, $91\mid f$ gives $h\ge\varphi(f)/3\ge\varphi(91)/3=24$ by the conductor formula, enough for both 19 and 31. In the split case, or if all four divide $f$, use the leading coefficients in the table below.
\begin{center}\small
\begin{tabular}{@{}ll@{}}\toprule
First prime coprime to $f$, if split & Five leading coefficients \\ \midrule
7 & $1,7,49,343,2401$ \\
13 & $1,13,49,169,2197$ \\
19 & $1,19,49,169,361$ \\
31 & $1,31,49,169,361$ \\
All four divide $f$ & $1,49,169,361,961$ \\ \bottomrule
\end{tabular}
\end{center}
Each row satisfies $4a_{\max}^2\le A$ and (2) at $D=A$ for every adjacent pair; both tests persist for larger $D$. This gives five separated complex conjugates and completes the large-discriminant proof.

\clearpage
\heading{The remaining finite interval.} We checked all $D\equiv0,3\pmod4$ with $3\le D<A$. The following disjoint counts are the output of \texttt{five/five\_sieve.cpp}.
\begin{center}\small
\begin{tabular}{@{}lr@{}}\toprule
First applicable certificate or reason for discharge & Discriminants \\ \midrule
Nonsplit at 7 or 13 away from the conductor & $9{,}109{,}793$ \\
Five complex moduli separated using (2) & $2{,}419{,}666$ \\
Complete enumeration gives $h<6$ & $25$ \\
Nonsplit at a larger prime and $h>4 S_p$ & $51$ \\
Individual class-polynomial test needed & $66$ \\ \midrule
Total & $11{,}529{,}601$ \\ \bottomrule
\end{tabular}
\end{center}
For odd $p$, the test for a conductor prime is $p\mid f$ iff $p^2\mid D$, since the odd part of $d_K$ is squarefree. The first pass discharges orders satisfying $h\ge6$ whenever the indicated congruence holds. The second searches reduced primitive forms with $a\le160$ for a chain passing (2). On survivors we count all reduced primitive forms up to $a\le\sqrt{D/3}$, counting both signs of $b$ except on the fundamental-domain boundary. For the fourth row we use $p=19,31,37,43,61,67,73,79,97,103$.

The $66$ residual polynomials, with $D\le4563$, have degrees from 6 through 12. A single disk certificate eliminates 62 of them for five terms. The target in (1) is $t=5$ regardless of the degree of the class polynomial. The archive includes an additional check on all orders of class number six, seven or eight using Watkins's fundamental-discriminant bounds. A five-root disk exists for 97 of 101 sextic orders, all 38 septic orders and all 208 octic orders. Including the residual list, this gives 367 disk witnesses among 371 distinct polynomials. Some certificate tables use larger depth thresholds than (1); they imply the stated test. The four possible exceptions are the sextics in the theorem. They each exclude \emph{four} terms directly: at $-87,-348$ use $(p,v,r)=(19,0,7)$ and at $-432,-675$ use $(11,0,1)$. The coefficient valuations of $H(p^v(r+Z))$ are $(4,3,2,1,0,0,0)$ in all four tests. Together with the independent quartic and quintic census, this also proves Theorem~\ref{th:four} without the earlier GRH cutoff or its $10^{11}$ sieve. No five-term witness is asserted for any of the five candidates.

\heading{Why another prime alone will not settle the five cases.}
For each remaining order, no five conjugates lie in a common multiplicative disk of depth strictly greater than $\kappa(p,5)$ at \emph{any} rational prime, even allowing algebraic centers. Five such roots would give 20 ordered pairs with $v_p(\alpha_j/\alpha_i-1)>\kappa(p,5)$. Such a pair implies $p\mid\mathrm{disc}(H)$, since the $\alpha_i$ are integral and $v_p(\alpha_j-\alpha_i)>0$. The roots of
\[
   Y^{-h}\operatorname{Res}_X\!\big(H(X),H((1+Y)X)\big)
\]
are the required ratios minus one, with the $h$ diagonal pairs removed. The bundled integer checker factors $|\mathrm{disc}(H)|$ by trial division and uses the least minimizing index at every candidate prime. In discriminant order $-87,-131,-348,-432,-675$, the maximum numbers of pairs are respectively $14,12,14,12,14$. This exhausts the particular sufficient condition (1) across all primes and centers. The width argument later in the paper uses four clustered units and an additional root of positive valuation.

\clearpage
\section{The power of the discriminant}\label{sec:growth}
We prove Theorem~\ref{th:growth}. For $p\ge3$ unramified in the splitting field at the chosen place, roots sharing a nonzero residue lie in a multiplicative disk of depth at least one. By Lemma~\ref{lem:disk} with $r=1$ and $v_p(k!)\le k/(p-1)$, a $t$-term equation has at most
\[
 \left\lfloor\frac{p-1}{p-2}(t-1)\right\rfloor
 \quad\hbox{distinct roots in each such disk}.                 \tag{3}
\]
This holds even when $t$ greatly exceeds $p$.

If $p\nmid D$ is inert in $K$, the ring class field is unramified above $p$: it is unramified over $K$ away from $f$, and $p$ is unramified in $K$. It contains all $h$ roots. After choosing a place, all their reductions are supersingular by Deuring's reduction theorem (potential good reduction suffices). The number of supersingular $j$-values over $\overline{\F}_p$ is $p/12+O(1)$, uniformly in $p$. Consequently if at least $h/2$ roots are units at this place, (3) gives $s(-D)\gg h/p$ \cite{cm,deuring}.

\heading{Selecting a small inert prime.} First suppose $K\ne\Q(\sqrt{-3})$. On the units modulo $q=6D$, the quadratic characters of $K$ and $\Q(\sqrt{-3})$ are distinct and nontrivial: they induce different primitive characters. The simultaneous conditions $\chi_K(p)=-1$, $\chi_{-3}(p)=1$ specify a coset of index four. Theorem 1.4 of Lamzouri, Li and Soundararajan \cite{lls} yields $p\ll (\log(2D))^2$ on Dirichlet GRH. Small moduli are absorbed in the absolute constant. This prime is coprime to $6D$ and $p\equiv1\pmod3$. The elliptic curve with $j=0$ has ordinary reduction at $p$, since $p$ splits in its CM field $\Q(\sqrt{-3})$. Hence every conjugate in our inert packet is a unit. The bound $s(-D)\gg h/(\log(2D))^2$ follows.

\heading{The exceptional CM field.} Here $D=3f^2$, and $f=1$ is immaterial. The complex tori of the principal lattices $\mathcal O_f\subset\mathcal O_K$ are linked by a cyclic isogeny of degree $f$; the latter has $j=0$. Pazuki's Theorem 1.1 gives $h_{\rm Weil}(\alpha)\le9.204+12\log f$ \cite{pazuki}. As $H_{-D}(0)\ne0$ is integral,
\[
  \log |H_{-D}(0)| \le h\,h_{\rm Weil}(\alpha).
\]
At an inert $p\nmid D$, each nonunit root contributes at least one to $v_p(H_{-D}(0))$, by unramifiedness. The total $\log p$ of such primes with more than half the roots nonunits is therefore at most $2h_{\rm Weil}(\alpha)$. Primes dividing $6f$ cost at most $\log(6f)$. On the other hand the primes $p\equiv2\pmod3$ up to $X$ have log-weight asymptotic to $X/2$, by the prime number theorem for this fixed progression. Taking $X=C\log(2f)$ for a sufficiently large absolute $C$ leaves an eligible prime with at least half the roots units. This actually proves $s(-3f^2)\gg h/\log(2f)$ without GRH.

Finally $h(-D)=D^{1/2+o(1)}$ for all orders, by the imaginary quadratic class-number formula, Siegel's lower bound and the conductor formula (spelled out on the next page). Together with $s\le h+1$, the uniform lower bound proves the asserted logarithmic limit.
\clearpage
\section{Two estimates independent of the zero theorem}
\heading{All discriminants.} For the order of conductor $f$ the class-number formula is
\[
 h(-D)=\frac{h(d_K)f}{[\mathcal O_K^\times:\mathcal O_f^\times]}
              \prod_{p\mid f}\left(1-\frac{\chi_{d_K}(p)}p\right).
\]
The unit index is at most three. The product is between $\varphi(f)/f$ and $\prod_{p\mid f}(1+1/p)$, both $f^{o(1)}$ as $f$ grows. Siegel's bound $h(d_K)\gg_\eta |d_K|^{1/2-\eta}$ and the elementary upper bound $h(d_K)\ll\sqrt{|d_K|}\log(2|d_K|)$ yield $h(-D)=D^{1/2+o(1)}$ uniformly. Enge's Theorem 1.2 gives $\log |H_{-D}(0)|\ll h+\sqrt D\,\log^2(2D)$ \cite{enge}. Discarding $D=3$, it follows that the number of inert $p\nmid D$ with more than $h/2$ nonunit roots is $D^{o(1)}$: each such prime costs at least $\frac12 h\log2$ from this same norm.

Apply Pollack's Theorem 1.1 \cite{pollack} to the Kronecker character of discriminant $-D$, including its imprimitive Euler factors. For each fixed $\delta>0$ there are $>D^{\kappa(\delta)}$ primes $p\le D^{1/(4\sqrt e)+\delta}$ at which the character takes value $-1$, for sufficiently large $D$. Every such prime is inert and coprime to $D$. After dropping 2, 3 and the $D^{o(1)}$ primes just counted, at least one remains. Apply (3) and $h=D^{1/2+o(1)}$, choosing $\delta<\varepsilon$ and absorbing the class-number loss into the difference. This proves
\[
      s(-D)\gg_\varepsilon D^{\,1/2-1/(4\sqrt e)-\varepsilon}.
\]
The exponent before $\varepsilon$ is approximately $0.348367$. Siegel's lower bound makes the constant in this proof ineffective.

\heading{Almost the largest possible support in two congruence families.}
At $p=7$ or $13$ there is exactly one supersingular $j$-value; it is nonzero since $p\equiv1\pmod3$ \cite{deuring}. If $(\Delta/p)=-1$, all $h$ conjugates lie in one unit disk of depth one. Let $s_p(k)$ denote the sum of the base-$p$ digits. With $t=s(\Delta)$, the proof of Lemma~\ref{lem:disk} gives some $k\ge h$ and $0\le i<t$ with
\[
 (p-2)k+s_p(k)\le(p-2)i+s_p(i),
\]
since $(p-1)(n-v_p(n!))=(p-2)n+s_p(n)$. Thus
\[
 h\le t-1+\frac{p-1}{p-2}(1+\log_p t).
\]
As $t\le h+1$, the claimed additive $O(\log(2h))$ loss follows. Order discriminants are all negative integers $\equiv0,1\pmod4$. The nonresidue conditions at 7 and 13 occur with proportions $3/7$ and $6/13$ independently, by the Chinese remainder theorem. Their union therefore has relative density $1-(4/7)(7/13)=9/13$. This gives a concrete unconditional near-equality even for the minimum over \emph{all} multiples of $H_\Delta$.
\clearpage
\section{An infinite family with exact minimum support}
The logarithmic loss can disappear entirely. This is an application of the same single-disk root count at a fixed prime; all its inputs are independent of RH.

\begin{theorem}\label{th:extreme}
Let $k\ge2$, $n=3\cdot7^{k-1}$ and $\Delta=-4\cdot 2^{2n}$. Then
\[
             s(\Delta)=h(\Delta)+1=2^{n-1}+1 .
\]
In particular every nonzero rational polynomial vanishing at a singular modulus of discriminant $-2^{44}$ has at least $1{,}048{,}577$ terms. The minimal polynomial attains that number.
\end{theorem}
\begin{proof}
Here $K=\Q(i)$ and $f=2^n$. The conductor formula gives $h=2^{n-1}$ because the prime 2 ramifies in $K$ and the unit index is two. The prime 7 is inert and prime to the conductor. All $h$ distinct roots lie in one multiplicative unit disk of depth one.

Write $W(m)=m-v_7(m!)=(5m+s_7(m))/6$. We prove that $W(i)<W(h)$ for $i<h$ and $W(i)>W(h)$ for $i>h$. It then follows that $B_7(h,1)<h$. A polynomial with at most $h$ terms cannot have the $h$ required zeros, by Lemma~\ref{lem:disk}.

Put $Q=7^k$. Since $2^3=1+7$, successive seventh powers show that $2h=2^n\equiv1\pmod Q$. For $2m+1<Q$ the valuations of $h+j$ equal those of $1+2j$ for $1\le j\le m$, as both numbers after multiplying the former by 2 agree modulo $Q$ and $0<1+2j<Q$. Similarly $v_7(h-j)=v_7(1-2j)$ for $0\le j<m$. Each of the corresponding two sums of valuations is at most $v_7((2m+1)!)$; in the second sum the only possible repeated absolute value is 1. For $m\ge1$ this bound is at most $(2m+1)/6<m$. Subtracting the factorial valuations proves the required strict inequalities at both offsets $m$.

For $m\ge(Q-1)/2$, use $s_7(i)\le6(1+\log_7 h)$ for $0\le i\le h$. Both $W(h+m)-W(h)$ and $W(h)-W(h-m)$, the latter for $m\le h$, are at least $(5m-6(1+\log_7 h))/6$. Since $\log_7 h<n/2=3Q/14$ and $5(Q-1)/2>6+9Q/7$ for $Q\ge49$, this is positive. The two ranges cover every integer $m\ge1$. The upper bound $s\le h+1$ completes the proof.
\end{proof}
The case $k=2$ has $n=21$ and $h=2^{20}$. The calculation proves that every coefficient of this million-degree class polynomial is nonzero. Multiplying it by a polynomial of any degree can never reduce its number of terms.

\clearpage
\section{The first remaining quintic}
At $\Delta=-131$, $5\le s(-131)\le6$. Two independent tests restrict a hypothetical five-term relation. Divide by its lowest monomial, so its exponents are $0=e_0<e_1<e_2<e_3<e_4$.
\begin{enumerate}\setlength{\itemsep}{.15em}
\item \emph{Bounded width.} For $e_4\le2048$ the five powers are independent over $\Q$. The checker reduces powers of $X$ modulo $H_{-131}$ over $\F_{2147483647}$ and drops the constant coordinate (the column of $e_0$). For each of the first two resulting four-vectors, fix their span, project the subsequent vectors into its two-dimensional quotient and compare projective directions. A collision is resolved by computing the rank at $1000000007$. These are prime moduli, so full rank at either proves rational independence. There are $\binom{2048}{4}=730{,}862{,}190{,}080$ normalized supports. All pass; 327 require the second prime.
\item \emph{Arbitrary width.} Some pair $1\le j<i\le4$ must satisfy $17^{e_j}\mid(e_i-e_j)$. In particular $e_4\ge e_1+17^{e_1}$.
\end{enumerate}
To prove the second test, the coefficient valuations at $(17,0,8)$ on page~\pageref{sec:cert} give four unit roots of the form $u_\ell=8\exp x_\ell$, with $v_{17}(x_\ell)\ge1$ by the Newton polygon, and the fifth root $w$ has valuation three since $v_{17}(H(0))=3$. For any four distinct exponents $b_j$ the determinant $\det(u_\ell^{b_j})$ has valuation
\[
 C+\sum_{i<j}v_{17}(b_j-b_i),
\]
where $C$ is independent of the exponents. Indeed expand $\det(\exp(x_\ell b_j))$ by Cauchy--Binet. For indices $0\le m_0<\cdots<m_3$, factor the Vandermondes in the $x$'s and $b$'s. Each remaining quotient of alternating monomials is an integral homogeneous symmetric polynomial of degree $d=\sum(m_i-i)$ (alternating determinant divisibility). Its $x$-evaluation has valuation at least $d$ and its $b$-evaluation has nonnegative valuation. The leading indices $m_i=i$ give $1/\prod_{i=0}^3i!$, a unit. For each $m_i>i$, $v_{17}(m_i!/i!)<m_i-i$, by the factorial bound on page~\pageref{lem:disk}; the remaining normalized terms have positive valuation. The exponential expansions converge, so the first term determines the valuation. The column factors $8^{b_j}$ are units.

Expand the full five-row determinant along $w$. If it vanishes, some term $j>0$ has valuation no larger than term 0. Canceling the common factors gives
\[
  3e_j\le\sum_{1\le i\le4,\ i\ne j}
       \big(v_{17}(e_j-e_i)-v_{17}(e_i)\big).
\]
One of the three summands is at least $e_j$. Its index exceeds $j$ because $17^{e_j}>e_j$. This proves the divisibility test. It also extends the exact computation to width $32{,}768$: only $e_1=1,2,3$ need to be considered there. Of $48{,}029{,}600{,}424{,}386{,}560$ total supports, the first-exponent filter leaves $17{,}587{,}354{,}664{,}945$; 8,275 modular collisions all pass the second check. Among supports of width $\le N$, at most $6N^3\sum_{m\ge1}17^{-m}=3N^3/8$ pass it: choose the pair, its smaller exponent $m$, the larger one in its congruence class and the two other exponents. Whether even one support yields a relation at unrestricted width remains open here.

\clearpage
\section{No five-term equation of width at most 65,536}
The \emph{width} of a polynomial is the difference between its largest and smallest exponents with nonzero coefficients. For a nonzero singular modulus we may divide the equation by its lowest power of $X$ and write $0=e_0<e_1<e_2<e_3<e_4$.

\begin{corollary}\label{cor:width}
Every singular modulus of degree at least five requires at least six terms in a rational equation of width at most $65{,}536$.
\end{corollary}
\begin{proof}
By Theorems~\ref{th:four} and~\ref{th:five} only five discriminants, and exactly five terms, require a test. At each discriminant the table gives four roots $u_1,\ldots,u_4$ with $v_p(u_i/r-1)>3/(p-1)$ and an additional root $w$ with $v_p(w)=v>0$. The checker verifies both assertions from the class polynomial. For the first, apply the least-minimizer test after translation; for the second, verify that the minimum of $v_p(A_i)+iv$ occurs at two or more indices.
\begin{center}\small
\begin{tabular}{@{}rrrrrr@{}}\toprule
$\Delta$ & $p$ & $r$ & $v$ & Pairs tested & Second-prime rechecks \\ \midrule
$-87$  & 23 & 19 & 3   & 2,977 & 2,983 \\
$-131$ & 17 &  8 & 3   & 4,094 & 3,994 \\
$-348$ & 29 &  2 & $3/2$ & 4,676 & 4,829 \\
$-432$ & 23 & 19 & 3   & 2,977 & 2,980 \\
$-675$ & 17 &  8 & 3   & 4,094 & 4,191 \\ \bottomrule
\end{tabular}
\end{center}
The four-row cofactor calculation from the preceding section still holds. Indeed every $x_\ell=\log(u_\ell/r)$ has valuation strictly greater than $3/(p-1)$. For $0\le i\le3$ and $m>i$, the factorial estimates proving (1) bound $v_p(m!/i!)/(m-i)$ by $1/(p-3)\le 3/(p-1)$. All nonleading Schur terms have strictly positive valuation after the two Vandermonde factors are removed. Expanding along $w$ now gives a pair $1\le j<i\le4$ with
\[
          p^{\lceil ve_j/3\rceil}\mid e_i-e_j .
\]
The larger exponent is indeed $e_i$, since $p^{\lceil v m/3\rceil}>m$ for every positive integer $m$ in these five cases.

For each admissible pair of exponents in $[1,65536]$, reduce all powers modulo the monic class polynomial over the two prime fields from the preceding section. Quotient out the constant column. For degree six project the remaining five coefficients $(a_1,\ldots,a_5)$ to $(a_1+a_5,\ldots,a_4+a_5)$; for degree five keep all four. A full-rank image proves independence. In the two-dimensional quotient by the fixed pair, distinct projective directions for the other two columns prove full rank at the first prime. All coincident directions are rechecked at the second prime. Every recorded recheck gives full rank. This verifies the claimed width for all five possible exceptions.
\end{proof}

\clearpage
\section{Verification and further questions}
\heading{Reproduction.} From the extracted directory, run \texttt{python3 verify.py}. The driver verifies the new certificate rows, regenerates the exact class polynomials, and replays the analytic endpoint arithmetic and both earlier discriminant sieves. Python 3 and a C++17 compiler suffice. The LaTeX manuscript builds with two runs of \texttt{pdflatex report.tex}. All reported finite exclusions are exact. An exponent search, when reported separately, certifies only its stated width. The local disk argument for Theorem~\ref{th:four} has no width bound.

\heading{Three remaining questions.} Does $s(-131)$ equal five or six? Can the four remaining sextic discriminants be settled for five terms? Finally, is the ratio $s(\Delta)/h(\Delta)$ bounded away from zero for all orders? The uniform estimate proves lower bounds with a logarithmic loss; the fixed-prime argument gives a ratio tending to one on the density-$9/13$ family. An answer to the last question would need control over discriminants outside that family as well.

\small
\begin{thebibliography}{12}
\setlength{\itemsep}{.05em}\setlength{\parskip}{0pt}
\bibitem{lean} L. Sernau and collaborators, \href{https://github.com/lsernau/grh-lean/tree/72fead102f663d5c7a407964a13dc0cf9d977207}{\emph{GRH for Dirichlet L-functions, in Lean 4}}. Readable edition, \repo{GUIDE.md}{guide}; based on the Alp\"oge--Claude proof.
\bibitem{audit} A. Sutherland, \href{https://math.mit.edu/~drew/shortgrh/}{\emph{Verification files for the shorter Lean proof of RH and GRH}}. Independent replay of commit \texttt{72fead10}.
\bibitem{prior} ProveIt, \href{https://github.com/VladimirReshetnikov/ProveIt/blob/13c30b3b677b691845466289bfdb792d11fefad5/docs/incoming/grh_cm_prime_applications.zip}{\emph{Sparse CM equations, prime-generating constants, and additive representations}} (incoming manuscript archive, 25 Nov. 2026). See \texttt{riffaut.tex}, \texttt{cm-finite.tex}, \texttt{cm\_small.tex}, \texttt{cm4075.tex}, \texttt{cm-conclusion.tex}.
\bibitem{parallel} \emph{Sparse equations for singular moduli and the sign of weighted prime sums}, parallel workspace manuscript, 25 Nov. 2026, archive \texttt{grh\_cm\_four\_terms\_and\_prime\_bias.zip}, sections on CM and growth.
\bibitem{riffaut} Y. Bilu, F. Luca and A. Pizarro-Madariaga, \href{https://arxiv.org/abs/2003.06547}{\emph{Trinomials, singular moduli and Riffaut's conjecture}}, especially \S\S2, 9.
\bibitem{few} M. Avenda\~no and T. Krick, \href{https://arxiv.org/abs/1008.4808}{\emph{Sharp bounds for the number of roots of univariate fewnomials}}, Prop. 3.8 and discussion of Lenstra's disk bounds.
\bibitem{gm} L. Greni\'e and G. Molteni, \href{https://arxiv.org/abs/1312.4463}{\emph{Explicit versions of the prime ideal theorem for Dedekind zeta functions under GRH}}, Cor. 1.2.
\bibitem{cm} D. A. Cox, \emph{Primes of the Form $x^2+ny^2$}, 2nd ed., Wiley, 2013, chs. 2--3. Ring class fields and the conductor class-number formula.
\bibitem{deuring} N. Elkies, K. Ono and T. Yang, \href{https://people.math.wisc.edu/~tonghaiyang/deuring.pdf}{\emph{Reduction of CM elliptic curves and modular function congruences}}, \S\S1--3. See also F. Oort, \href{https://webspace.science.uu.nl/~oort0109/EigArt-HistpRH-2016.pdf}{historical account of the supersingular count}.
\bibitem{lls} Y. Lamzouri, X. Li and K. Soundararajan, \href{https://arxiv.org/abs/1309.3595}{\emph{Conditional bounds for the least quadratic non-residue and related problems}}, Thm. 1.4.
\bibitem{pazuki} F. Pazuki, \href{https://arxiv.org/abs/1611.01094}{\emph{Modular invariants and isogenies}}, Thm. 1.1.
\bibitem{enge} A. Enge, \href{https://arxiv.org/abs/cs/0601104}{\emph{The complexity of class polynomial computation via floating point approximations}}, Thm. 1.2.
\bibitem{pollack} P. Pollack, \href{https://arxiv.org/abs/1508.05035}{\emph{Bounds for the first several prime character nonresidues}}, Thm. 1.1.
\bibitem{watkins} M. Watkins, \href{https://magma.maths.usyd.edu.au/~watkins/papers/smallcn.pdf}{\emph{Class numbers of imaginary quadratic fields}}, Table 4.
\bibitem{orders} \href{https://arxiv.org/html/2605.26020v1}{\emph{Boundary CM points and class groups of small exponent}}, Table A.5 (all order discriminants of class number five).
\end{thebibliography}
\end{document}
